% TOP_PROBLEM: {"id": 3156, "title": "Friendly partitions", "category_id": 3, "category": {"id": 3, "name": "graph_theory", "display_name": "Graph Theory", "description": "Problems involving graphs, networks, and their properties.", "slug": "graph-theory", "order_index": 3, "created_at": "2026-07-31T15:26:25.671Z"}, "status": "open", "problem_number": "OPG-37163", "set_id": 12, "statusQualification": "Both partition classes are required to be nonempty, as in the standard internal-partition problem; the source short definition leaves this implicit."}
% FIRST_POSTED: 2026-10-01
% ENTRY_KIND: statement_only
% UnsolvedMath ID 3156; source catalogue code OPG-37163
% !TeX program = lualatex
% Definition and sources only; this file is not an LLM solution attempt.
\documentclass[11pt,a4paper]{article}
\usepackage[margin=25mm]{geometry}
\usepackage{fontspec}
\setmainfont{lmroman10-regular.otf}[BoldFont=lmroman10-bold.otf,
 ItalicFont=lmroman10-italic.otf,BoldItalicFont=lmroman10-bolditalic.otf]
\usepackage{amsmath,amssymb,mathtools}
\usepackage{unicode-math}
\setmathfont{latinmodern-math.otf}
\AtBeginDocument{\let\setminus\smallsetminus}
\usepackage{xurl,hyperref}
\hypersetup{colorlinks=true,urlcolor=blue}
\setcounter{secnumdepth}{0}
\setlength{\parindent}{0pt}
\setlength{\parskip}{5pt}
\setlength{\emergencystretch}{3em}
\begin{document}
\section*{Friendly partitions}\label{3156}
{\small UnsolvedMath ID 3156 · OPG-37163 · Graph Theory · OpenGarden}
\subsection{Definitions and mathematical statement}
Let $G=(V,E)$ be a finite simple graph. For $v\in V$ and
$X\subseteq V$, put $d_X(v)=|N_G(v)\cap X|$. A friendly, or
internal, partition is a decomposition
$V=A\mathbin{\dot\cup}B$ into two nonempty sets such that
\[
 d_A(v)\geq d_B(v)\quad(v\in A),\qquad
 d_B(v)\geq d_A(v)\quad(v\in B).
\]
The requirement that both parts be nonempty excludes the trivial
partition putting every vertex on the same side. No equality of
the two part sizes is required.

For an $r$-regular graph, these inequalities are equivalent to
each vertex having at least $\lceil r/2\rceil$ neighbours in its
own part. In particular, each induced subgraph $G[A]$ and $G[B]$
must have minimum degree at least $\lceil r/2\rceil$.

The conjecture asks whether for every fixed integer $r\geq1$
there exists an integer $N(r)$ such that every finite simple
$r$-regular graph on at least $N(r)$ vertices has a friendly
partition. This is the meaning of all but finitely many graphs,
with exceptions counted up to isomorphism. The degree is fixed
before the order grows. The threshold must work for every graph
of that degree, not only almost every randomly sampled graph.
Disconnected graphs are included; the difficult instances are
the connected ones. The assertion is local at every vertex,
not just a bound on the total number of crossing edges.
\subsection{Short English statement}
For each fixed degree, do all sufficiently large regular graphs split into two nonempty parts so every vertex has at least half its neighbours on its own side?
\subsection{Sources}
\begin{itemize}
\item[S1] ULAM.AI and the UnsolvedMath Contributors, supplied problems.json, record 3156 (OPG-37163), OpenGarden collection. Catalogue identity and source transcription; CC BY 4.0.
\url{https://huggingface.co/datasets/ulamai/UnsolvedMath}.
\item[S2] Open Problem Garden, Friendly partitions. Original problem and discussion; the mathematical formulation below is expanded from the supplied source record.
\url{http://www.openproblemgarden.org/op/friendly_partitions}.
\item[S3] A. Ban and N. Linial, Internal Partitions of Regular Graphs (2013), definitions and fixed-degree formulation.
\url{https://arxiv.org/abs/1307.5246}.
\end{itemize}
\end{document}
